\documentclass{article}
\usepackage{mathtools} 
\usepackage{fontspec}
\usepackage[UTF8]{ctex}
\usepackage{amsthm}
\usepackage{mdframed}
\usepackage{xcolor}
\usepackage{amssymb}
\usepackage{amsmath}


% 定义新的带灰色背景的说明环境 zremark
\newmdtheoremenv[
  backgroundcolor=gray!10,
  % 边框与背景一致，边框线会消失
  linecolor=gray!10
]{zremark}{说明}

% 通用矩阵命令: \flexmatrix{矩阵名}{元素符号}{行数}{列数}
\newcommand{\flexmatrix}[4]{
  \[
  #1 = \begin{pmatrix}
    #2_{11}     & #2_{12}     & \cdots & #2_{1#4}   \\
    #2_{21}     & #2_{22}     & \cdots & #2_{2#4}   \\
    \vdots      & \vdots      & \ddots & \vdots     \\
    #2_{#31}    & #2_{#32}    & \cdots & #2_{#3#4}
  \end{pmatrix}
  \]
}

% 简化版命令（默认矩阵名为A，元素符号为a）: \quickmatrix{行数}{列数}
\newcommand{\quickmatrix}[2]{\flexmatrix{A}{a}{#1}{#2}}

\begin{document}
\title{习题 12.2}
\author{张志聪}
\maketitle

\section*{1}

\begin{itemize}
  \item (4)

        当$N \geq e^2$，则$n > N$有
        \begin{align*}
          (\ln n)^n > 2^n
        \end{align*}
        于是
        \begin{align*}
          \frac{1}{(\ln n)^n} < \frac{1}{2^n}
        \end{align*}
        $\sum \frac{1}{2^n}$收敛，故$\sum\limits_{n = 2}^{\infty} \frac{1}{(\ln n)^n}$收敛。

  \item (7)

        我们有$e^x - 1 \sim x \, (x \to 0)$，
        又有
        \begin{align*}
          \lim\limits_{x \to 0} \frac{a^x -1}{e^x - 1}
           & = \lim\limits_{x \to 0} \frac{a^x \ln a}{e^x}                       \\
           & = \lim\limits_{x \to 0} (\frac{a}{e})^x \ln a                       \\
           & = \lim\limits_{x \to 0} (\frac{a}{e})^x \lim\limits_{x \to 0} \ln a \\
           & = 1 \cdot \ln a                                                     \\
           & = \ln a
        \end{align*}
        于是可得
        \begin{align*}
          \lim\limits_{n \to \infty} \frac{a^{\frac{1}{n}} -1}{\frac{1}{n}}
           & = \lim\limits_{x \to 0} \frac{a^{x} -1}{x}       \\
           & = \lim\limits_{x \to 0} \frac{a^{x} -1}{e^x - 1} \\
           & = \ln a
        \end{align*}
        因为$\sum \frac{1}{n}$发散，故$\sum a^{\frac{1}{n}} -1 $发散。

  \item (8)
        \begin{align*}
          e^{\ln\, (\ln n)^{\ln n}}
           & = e^{\ln n \cdot \ln (\ln n)}          \\
           & = \left[e^{\ln n}\right]^{\ln (\ln n)} \\
           & = n^{\ln (\ln n)}
        \end{align*}
        因为$\lim \limits_{n \to \infty} \ln (\ln n) = +\infty$，
        于是存在$N > 0$，$n > N$有
        \begin{align*}
          \ln (\ln n) > 3
        \end{align*}
        所以，$n > N$，有
        \begin{align*}
          \frac{n}{e^{\ln\, (\ln n)^{\ln n}}} < \frac{n}{n^3} = \frac{1}{n^2}
        \end{align*}
        $\sum \frac{1}{n^2}$收敛，故$\sum \frac{n}{ (\ln n)^{\ln n}}$收敛。

  \item (9)

        泰勒展开
        \begin{align*}
          a^x    & = 1 + \ln a \cdot x + \frac{(\ln a)^2 x^2}{2!} + o(x^2) \, (x \to 0) \\
          a^{-x} & = 1 - \ln a \cdot x + \frac{(\ln a)^2 x^2}{2!} + o(x^2) \, (x \to 0)
        \end{align*}
        所以
        \begin{align*}
          a^x + a^{-x} - 2 & = (\ln a)^2 x^2 + o(x^2) \, (x \to 0)
        \end{align*}
        综上
        \begin{align*}
          \lim\limits_{n \to \infty} \frac{a^{\frac{1}{n}} + a^{-{\frac{1}{n}}} - 2}{\frac{1}{n^2}}
           & = \lim\limits_{x \to 0} \frac{a^x + a^{-x} - 2}{x^2}       \\
           & = \lim\limits_{x \to 0} \frac{(\ln a)^2 x^2 + o(x^2)}{x^2} \\
           & = (\ln a)^2
        \end{align*}
        $\sum \frac{1}{n^2}$收敛，故$\sum a^{\frac{1}{n}} + a^{-{\frac{1}{n}}} - 2$收敛。

  \item (10)
        \begin{align*}
          \lim \limits_{n \to \infty} 2 n sin \frac{1}{n}
           & = \lim \limits_{n \to \infty} 2\frac{sin \frac{1}{n}}{\frac{1}{n}} \\
           & = 2
        \end{align*}
        所以有保号性可知，存在$N > 0$， $n > N$有
        \begin{align*}
          2\frac{sin \frac{1}{n}}{\frac{1}{n}} > \frac{3}{2}
        \end{align*}
        于是，$n > N$，有
        \begin{align*}
          \frac{1}{n^{2 n sin \frac{1}{n}}} < \frac{1}{n^{3/2}}
        \end{align*}
        $\sum \frac{1}{n^{3/2}}$收敛，故$\sum \frac{1}{n^{2 n sin \frac{1}{n}}}$收敛。
\end{itemize}

\section*{2}

\begin{itemize}
  \item (4)
        \begin{align*}
          \lim\limits_{n \to \infty} \frac{u_{n + 1}}{u_n}
           & = \lim\limits_{n \to \infty} \frac{(n +1)!}{(n+1)^{n + 1}} \cdot \frac{n^n}{n!}                                    \\
           & = \lim\limits_{n \to \infty} \frac{(n +1) n^n}{(n+1)^{n + 1}}                                                      \\
           & = \lim\limits_{n \to \infty} \frac{(n +1) n^n}{(n +1)(n+1)^{n}}                                                    \\
           & = \lim\limits_{n \to \infty} \frac{n^n}{(n+1)^{n}}                                                                 \\
           & = \lim\limits_{n \to \infty} \left(\frac{n}{n+1}\right)^n                                                          \\
           & = \lim\limits_{n \to \infty} \left(1 - \frac{1}{n + 1}\right)^n                                                    \\
           & = \lim\limits_{n \to \infty} \left(1 + \frac{1}{-(n + 1)}\right)^n                                                 \\
           & = \lim\limits_{n \to \infty} \frac{\left(1 + \frac{1}{-(n + 1)}\right)^{n+1}}{\left(1 + \frac{1}{-(n + 1)}\right)}
        \end{align*}
        考虑
        \begin{align*}
          \lim\limits_{n \to \infty} \left(1 + \frac{1}{-(n + 1)}\right)^{n+1}
           & = \lim\limits_{n \to \infty} \frac{1}{\left(1 + \frac{1}{-(n + 1)}\right)^{-(n+1)}} \\
           & = \frac{1}{e}
        \end{align*}
        （使用了归结原理：$\lim\limits_{x \to -\infty} (1 + \frac{1}{x})^x = e$）。

        又因为
        \begin{align*}
          \lim\limits_{n \to \infty} 1 + \frac{1}{-(n + 1)} = 1
        \end{align*}
        所以
        \begin{align*}
          \lim\limits_{n \to \infty} \frac{\left(1 + \frac{1}{-(n + 1)}\right)^{n+1}}{\left(1 + \frac{1}{-(n + 1)}\right)}
           & = \frac{\frac{1}{e}}{1} \\
           & = \frac{1}{e} \\
           & < 1
        \end{align*}
        故收敛。
\end{itemize}


\end{document}